\documentclass{beamer}

% You can also use a 16:9 aspect ratio:
%\documentclass[aspectratio=169]{beamer}
\usetheme{TACC16}

% It's possible to move the footer to the right:
%\usetheme[rightfooter]{TACC16}

\begin{document}
\title[Lmod]{Lmod: A Modern Environment Module System}
\author{Robert McLay} 
\date{April 20, 2017} 

% page 1
\frame{\titlepage} 

% page 2
\begin{frame}{Outline}
  \begin{itemize}
    \item What are Environment Modules?
    \item What is Lmod?
    \item Modulefiles Architecture
    \item Advanced Topics
    \item Where to go for help.
  \end{itemize}
\end{frame}

% page 3
\begin{frame}{Main Takeaways}
  \begin{itemize}
    \item Latest version: https://github.com:TACC/Lmod.git
    \item Stable version: http://lmod.sf.net
    \item Documentation:  http://lmod.readthedocs.org
    \item Mailing List:   lmod-users@lists.sourceforge.net.
    \item Join here: https://lists.sourceforge.net/lists/listinfo/lmod-users
  \end{itemize}
\end{frame}

% page 4
\begin{frame}{What are Modules?}
  \begin{itemize}
    \item Modules are the way sites provide optional software: MPI,
      Boost, ...
    \item Modules add to PATH and set other env. vars for each package
    \item Modules can be unloaded.
    \item Modulefiles are in one file not one for each shell. (e.g. Compiler init scripts)
  \end{itemize}
\end{frame}

% page 5
\begin{frame}{What is Lmod?}
  \begin{itemize}
    \item A modern replacement for a tried and true concept.
    \item The guiding principal: ``Make life easier without getting in
      the way.''
    \item Reads both TCL and Lua modulefiles
  \end{itemize}
\end{frame}

% page 6
\begin{frame}{Fundamental Issues}
  \begin{itemize}
    \item Software Packages are created and updated all the time.
    \item Some Users need new versions for new features and bug fixes.
    \item Other Users need older versions for stability and continuity.
    \item No system can support all versions of all packages.
    \item User programs using pre-built C++ \& Fortran libraries must link with the same compiler.
    \item Similarly, MPI Applications must build and link with same
      MPI/Compiler pairing when using pre-built MPI libraries.
  \end{itemize}
\end{frame}

% page 7
\begin{frame}[fragile]
    \frametitle{Example of Lmod: Environment Modules (I)}
    {\tiny
\begin{semiverbatim}
\$ {\color{blue} module avail}
------------------ /opt/apps/modulefiles/MPI/intel/12.0/mpich2/1.4 ------------------
  petsc/3.1 (default)    petsc/3.1-debug    pmetis/4.0    tau/2.20.3

------------------- /opt/apps/modulefiles/Compiler/intel/12.0 -----------------------
  boost/1.45.0              gotoblas2/1.13          openmpi/1.4.3
  boost/1.46.0              mpich2/1.3.2            openmpi/1.5.1
  boost/1.46.1 (default)    mpich2/1.4 (default)    openmpi/1.5.3 (default)

-------------------------- /opt/apps/modulefiles/Core -------------------------------
  StdEnv               intel/11.1               papi/4.1.4
  admin/admin-1.0      intel/12.0 (default)     scite/2.28
  ddt/ddt              lmod/lmod                tex/2010
  dmalloc/dmalloc      local/local (default)    unix/unix (default)
  fdepend/1.2          mkl/mkl                  visit/visit
  gcc/4.4              noweb/2.11b
  gcc/4.5 (default)
\end{semiverbatim}
    }
\end{frame}

% page 8
\begin{frame}[fragile]
    \frametitle{Example of Lmod: Environment Modules (II)}
    {\tiny
\begin{semiverbatim}
{\color{blue}\$ module list}
Currently Loaded Modules:
  1) StdEnv  2) gcc/4.5  3) mpich2/1.4  4) petsc/3.1
{\color{blue}\$ module unload gcc}
Inactive Modules:
  1) mpich2  2) petsc
{\color{blue}\$ module list}
Currently Loaded Modules:
  1) StdEnv
Inactive Modules:
  1) mpich2  2) petsc
{\color{blue}\$ module load intel}
Activating Modules:
  1) mpich2  2) petsc
{\color{blue}\$ module swap intel gcc}
Due to MODULEPATH changes the follow modules have been reloaded:
  1) mpich2  2) petsc
\end{semiverbatim}
    }
\end{frame}

% page 9
\begin{frame}{Why You Might Want To Use Lmod?}
  \begin{itemize}
    \item Same \texttt{module} commands as in Tmod (modules.sf.net)
    \item Active Development;  Frequent Releases; Bug fixes.
    \item Vibrant Community
    \item It is used from Norway to Israel to New Zealand from Stanford to MIT to NASA
    \item Enjoy many capabilities w/o changing a single module file
    \item Debian, Fedora and Brew packages available
    \item Many more advantages when you're ready
    \item It is what we and many sites around the world use every day!
  \end{itemize}
\end{frame}

% page 10
\begin{frame}{Lmod Features}
  \begin{itemize}
    \item Reads both TCL and Lua modulefiles (mixing works fine)
    \item One name rule.
    \item Support a Software Hierarchy
    \item Fast \texttt{module avail} via optional spider cache 
    \item Properties (gpu, mic)
    \item Semantic Versioning:  5.6 is older than 5.10
    \item family(``compiler''), family(``MPI'') support
    \item Optional Tracking: What modules are used?
    \item Many other features: ml, collections, hooks, nag, ...
  \end{itemize}
\end{frame}

% page 11
\begin{frame}{Tmod vs. Lmod}
  \begin{itemize}
    \item Tmod (TCL/C) is in maintenance mode, Lmod active
    \item Some work on Tmod (TCL) recently
    \item Lmod has many more features
    \item Tmod: \texttt{module load gcc/5.3 gcc/6.0} works
    \item Lmod has the ``One Name Rule''
    \item Lmod is close to Tmod, but not the same.
  \end{itemize}
\end{frame}

% page 12
\begin{frame}{Safety Features of Lmod (I)}
  \begin{itemize}
    \item Users can only load one version of a package.
    \item ``\texttt{module load xyz/2.1}'' loads xyz version 2.1
    \item ``\texttt{module load xyz/2.2}'' unloads 2.1, loads 2.2
    \item No way to force Lmod to load two or more version with the
      same ``name''!
  \end{itemize}
\end{frame}

% page 13
\begin{frame}{Safety Features of Lmod (II)}
  \begin{itemize}
    \item Lmod adds a new command in modulefiles: \texttt{family("}\emph{name}\texttt{")}
    \item All of our compiler modules have \texttt{family("compiler")}
    \item All of our MPI modules have \texttt{family("MPI")}
    \item Users can only load one compiler or MPI at a time
    \item Power users can get around this restriction.
  \end{itemize}
\end{frame}

%page 14
\begin{frame}{Module Architecture Design}
  \begin{itemize}
    \item Lua or TCL modulefiles?
    \item Optional software: shared or local?
    \item Flat or Hierarchical Module Layout?
    \item Naming conventions? (N/V, C/N/V, N/V/V?)
    \item A standard set of modules?
    \item Keep software for life of cluster or not?
    \item Problems with Bash
  \end{itemize}
\end{frame}

%page 15
\begin{frame}{Shared Disk versus Local Install}
  \begin{itemize}
    \item Local install: Fast, small 
    \item Shared Disk: Big, slow
    \item TACC: local install (mostly)
    \item Limited to 3 version max!
    \item Shared Disk: Keep software for life of system?
  \end{itemize}
\end{frame}

%page 16
\begin{frame}{Modulefile Choices}
  \begin{itemize}
    \item Flat Naming Scheme?
    \item Hierarchical Naming Scheme?
  \end{itemize}
\end{frame}

%page 17
\begin{frame}{Flat Naming Scheme: PETSc}
  PETSc is a parallel iterative solver package:
  \begin{itemize}
    \item Compilers: GCC 6.3, Intel 17.0
    \item MPI Implementations: MVAPICH2 2.1, IMPI 17.0
    \item MPI Solver package: PETSc 4.1
    \item 4 versions of PETSc: 2 Compilers $\times$ 2 MPI
  \end{itemize}
\end{frame}

%page 18
\begin{frame}{Flat: PETSc }
  \begin{enumerate}
  \item \texttt{PETSc/4.1-mvapich2-2.1-gcc-6.3}
  \item \texttt{PETSc/4.1-mvapich2-2.1-intel-17.0}
  \item \texttt{PETSc/4.1-openmpi-1.8-gcc-6.3}
  \item \texttt{PETSc/4.1-openmpi-1.8-intel-17.0}
  \end{enumerate}
\end{frame}

%page 19
\begin{frame}{Problems w/ Flat naming scheme}
  \begin{itemize}
    \item Users have to load modules:
      \begin{itemize}
        \item ``intel/17.0''
        \item ``mvapich2/2.1-intel-17.0''
        \item ``PETSc/4.1-mvapich2-2.1-intel-17.0''
        \item Changing compilers means unloading all three modules
        \item Reloading new compiler, MPI, PETSc modules.
        \item Not loading correct modules $\Rightarrow$ Mysterious Failures!
        \item Onus of package compatibility on users!
        \item Or extremely complicated modulefiles!
      \end{itemize}
  \end{itemize}
\end{frame}

%page 20
\begin{frame}{Extremely complicated modulefiles}
  \begin{itemize}
    \item Protect users via conflicts and/or prereqs
    \item The problem is that they are fragile
    \item What happens with a new compiler or MPI stack?
  \end{itemize}
\end{frame}

%page 21
\begin{frame}{Hierarchical Naming Schemes}
  \begin{itemize}
    \item Store modules under one tree: \texttt{/opt/apps/modulefiles}
    \item One strategy is to use sub-directories:
      \begin{itemize}
        \item Core: Regular packages: apps, compilers, git
        \item Compiler: Packages that depend on compiler: boost, MPI
        \item MPI: Packages that depend on MPI/Compiler: PETSc, FFTW3
      \end{itemize}
  \end{itemize}
\end{frame}

%page 22
\begin{frame}{\texttt{MODULEPATH}}
  \begin{itemize}
    \item \texttt{MODULEPATH} is a colon separated list of directories
      containing directories and module files.
    \item No modulefiles loaded $\Rightarrow$ users can only load core modules.
    \item Loading a compiler module adds to the \texttt{MODULEPATH}:
      \begin{itemize}
        \item Users can load compiler dependent modules.
        \item This includes MPI implementations modules.
      \end{itemize}
    \item Loading an MPI module adds to the \texttt{MODULEPATH}:
      \begin{itemize}
        \item Users can load MPI libraries that match the MPI/compiler pairing
      \end{itemize}
  \end{itemize}
\end{frame}


%page 23
\begin{frame}{Hierarchical Examples: Core}
  \begin{itemize}
    \item Generic:
      \begin{itemize}
        \item Package: \texttt{/opt/apps/}\emph{package/version}
        \item M: {\color{blue}/opt/apps/modulefiles}
        \item Modulefile: \texttt{{\color{blue}\$M}/Core/}\emph{package/version}
      \end{itemize}
    \item Git 1.8
      \begin{itemize}
        \item Package: \texttt{/opt/apps/git/1.8}
        \item Modulefile: \texttt{{\color{blue}\$M}/Core/git/1.8}
      \end{itemize}
    \item Intel compilers 17.0
      \begin{itemize}
        \item Package: \texttt{/opt/apps/intel/17.0}
        \item Modulefile: \texttt{{\color{blue}\$M}/Core/intel/17.0}
        \item Modulefile adds \texttt{{\color{blue}\$M}/Compiler/intel/17.0} to \texttt{MODULEPATH}
      \end{itemize}
  \end{itemize}
\end{frame}

%page 24
\begin{frame}{Hierarchical Examples: Compiler Dependent}
  \begin{itemize}
    \item Generic:
      \begin{itemize}
        \item Package: /opt/apps/\emph{compiler-version/package/version}
        \item M: {\color{blue}/opt/apps/modulefiles}
        \item Modulefile: \texttt{{\color{blue}\$M}/Compiler/}\emph{compiler/version/package/version}
      \end{itemize}
    \item Openmpi 1.8 with gcc 6.3
      \begin{itemize}
        \item Package: \texttt{/opt/apps/gcc-6\_3/openmpi/1.8}
        \item Modulefile: \texttt{{\color{blue}\$M}/Compiler/gcc/6.3/openmpi/1.8}
        \item Modulefile adds \texttt{{\color{blue}\$M}/MPI/gcc/6.3/openmpi/1.8}
          to \texttt{MODULEPATH} \\
      \end{itemize}
    \item Openmpi 1.8 with intel 17.0
      \begin{itemize}
        \item Package: \texttt{/opt/apps/intel-17\_0/openmpi/1.8}
        \item Modulefile: \texttt{{\color{blue}\$M}/Compiler/intel/17.0/openmpi/1.8}
        \item Modulefile adds \texttt{\$M/MPI/intel/17.0/openmpi/1.8}
          to \texttt{MODULEPATH}
      \end{itemize}
  \end{itemize}
\end{frame}

%page 25
\begin{frame}{Hierarchical Examples: MPI/Compiler Dependent}
  \begin{itemize}
    \item PETSc 4.1 (1/4)
      \begin{itemize}
        \item Package: \texttt{/opt/apps/intel-17\_0/openmpi-1\_8/petsc/4.1}
        \item Modulefile: \texttt{{\color{blue}\$M}/MPI/intel/17.0/openmpi/1.8/petsc/4.1}
      \end{itemize}
    \item PETSc 4.1 (2/4)
      \begin{itemize}
        \item Package: \texttt{/opt/apps/gcc-4\_5/mvapich2-2\_1/petsc/4.1}
        \item Modulefile: \texttt{{\color{blue}\$M}/MPI/gcc/6.3/mvapich2/2.1/petsc/4.1}
      \end{itemize}
  \end{itemize}
\end{frame}


%page 26
\begin{frame}{Loading the correct module}
  \begin{itemize}
    \item User loads ``\texttt{intel/17.0}'' module
    \item Can only see/load compiler dependent packages that are built with
      intel 17.0 compiler.
    \item Can not see/load package built with other versions or other compilers.
    \item Similar loading ``\texttt{openmpi/1.8}'' module.
    \item Users can only load package that are built w/ intel 17.0 and
      openmpi 1.8 and no others.
  \end{itemize}
\end{frame}


%page 27
\begin{frame}{Modulefile contents}
  \begin{itemize}
    \item Be consistent! Find a convention and stick with it
    \item We define consistent variables in each module:
      $<$SITE\_NAME$>$\_$<$PKG\_Name$>$\_\{LIB,INC,BIN\}
      \begin{itemize}
        \item TACC\_HDF5\_INC
        \item TACC\_HDF5\_BIN
      \end{itemize}
  \end{itemize}
\end{frame}


%page 28
\begin{frame}[fragile]
    \frametitle{A Standard Set of Modules}
  \begin{itemize}
    \item Define a compiler and mpi stack
    \item /opt/apps/modulefiles/Core/StdEnv.lua
  {\small
    \begin{semiverbatim}
load("gcc","mvapich2")
    \end{semiverbatim}
}
  \end{itemize}
\end{frame}


%page 29
\begin{frame}[fragile]
    \frametitle{A Standard Set of Modules (II)}
  \begin{itemize}
    \item /etc/profile.d/z01\_StdEnv.sh
  {\small
    \begin{semiverbatim}
   if [ -z "\$\_\_Init\_Default\_Modules" ]; then
      export \_\_Init\_Default\_Modules=1;

      export LMOD\_SYSTEM\_DEFAULT\_MODULES="StdEnv"
      module --initial\_load --no\_redirect restore
   else
      module refresh
   fi
    \end{semiverbatim}
}
    \item Why is  ``module refresh'' there?
  \end{itemize}
\end{frame}


%page 30
\begin{frame}{User collections}
  \begin{itemize}
    \item Users may save a ``default'' collection of modules 
    \item module save 
    \item Users may save as many named collections as they like
    \item module save \emph{name}
    \item module restore \emph{name} will remove all modules then load
      the collection.
  \end{itemize}
\end{frame}

%page 31
\begin{frame}{Save/Restore}
  \begin{itemize}
    \item Users can setup their own initially loaded modules.
      \begin{itemize}
        \item Users simply load, unload and/or swap until happy.
        \item {\color{blue}\texttt{module save}} saves state in ``default''
        \item Our login scripts do: {\color{blue}\texttt{module restore}}
          which loads the user's default.
      \end{itemize}
    \item Users can create other collections by:
      \begin{itemize}
        \item {\color{blue}\texttt{module save}
            {\color{violet}name}} to save it.
        \item {\color{blue}\texttt{module restore}
            {\color{violet}name}} to retrieve it.
      \end{itemize}
  \end{itemize}
\end{frame}

%page 32
\begin{frame}{Module reset, restore}
  \begin{itemize}
    \item module reset $\rightarrow$ module purge; module load
      \$LMOD\_SYSTEM\_DEFAULT\_MODULES
    \item module restore looks for a default collection or loads the
      system default.
  \end{itemize}
\end{frame}

%%page 
%\begin{frame}[fragile]
%    \frametitle{}
%  \begin{itemize}
%    \item 
%  {\small
%    \begin{semiverbatim}
%
%    \end{semiverbatim}
%}
%    \item 
%  \end{itemize}
%\end{frame}

%page 33
\begin{frame}{Module Naming Convenions}
  \begin{itemize}
    \item N/V:   Name/Version (e.g. \texttt{bowtie/2.3})
    \item C/N/V: Category/Name/Version (e.g. \texttt{bio/bowtie/2.3})
    \item N/V/V: Name/Version/Version (e.g. \texttt{bowtie/64/2.3})
  \end{itemize}
\end{frame}

%page 34
\begin{frame}{Module Naming Conventions (II)}
  \begin{itemize}
    \item Try to stick with N/V if possible
    \item It's less typing
    \item C/N/V might be helpful to novice users
    \item But your obvious categories may not be obvious to your users
    \item Avoid N/V/V unless your users are experts 
    \item Or if you really need 64/32 bit libraries
  \end{itemize}
\end{frame}

%page 35
\begin{frame}{Bash Issues}
  \begin{itemize}
    \item Bash Startup is typically ``broken'' for non-login interactive shells
    \item Redhat, Centos, MacOS typically don't source /etc/bashrc on interactive shells
    \item MPI jobs start an interactive shell.
  \end{itemize}
\end{frame}

%page 36
\begin{frame}{Bash Issues (II)}
  \begin{itemize}
    \item Want module command to work in all shells.
    \item Want stacksize unlimited for MPI jobs
    \item We patched bash to force it to source /etc/tacc/bashrc
  \end{itemize}
\end{frame}

%page 37
\begin{frame}{Bash Repair Choices}
  \begin{itemize}
    \item Switch users to Z shell?
    \item patch bash (see Lmod docs)
    \item Expect all users to source /etc/bashrc in \textasciitilde/.bashrc
    \item Expect all users to start jobs with \#!/bin/bash -l
  \end{itemize}
\end{frame}

%page 38
\begin{frame}{Keeping Software for Life of Cluster or Not}
  \begin{itemize}
    \item It is possible with a shared disk approach
    \item You might want to hide older modules
  \end{itemize}
\end{frame}

%page 39
\begin{frame}{Hidden Modules}
  \begin{itemize}
    \item Modules can be hidden by adding a leading dot
    \item For example gcc/.6.3
    \item Lmod 7 also allows for hidden modules via MODULERC
    \item system MODULERC or \textasciitilde/.modulerc or file .modulerc
  \end{itemize}
\end{frame}

%page 40
\begin{frame}[fragile]
    \frametitle{Using System MODULERC to hide modules}
  \begin{itemize}
    \item In \$MODULERCFILE
  {\small
    \begin{semiverbatim}
        #%Module
        hide-version foo/3.2
    \end{semiverbatim}
}
  \end{itemize}
\end{frame}

%page 41
\begin{frame}{Deprecating Modules}
  \begin{itemize}
    \item Sometime you may have to remove a package (and its
      modulefile)
    \item Certainly with Local installs.
    \item How do you let users know?
  \end{itemize}
\end{frame}

%page 42
\begin{frame}{Deprecating Modules (II)}
  \begin{itemize}
    \item Assume Lmod is installed as /opt/apps/lmod/7.4.1/...
    \item Then /opt/apps/lmod/etc/admin.list is the ``nag'' file
    \item Or set \$LMOD\_ADMIN\_FILE to any file.
    \item Every time a particular modulefile is loaded the user gets a
      message.
    \item We have found that setting a nag message to be very effective!
  \end{itemize}
\end{frame}

%page 43
\begin{frame}{Tracking Module Usage}
  \begin{itemize}
    \item Lmod makes it easy to track module usage.
    \item Lmod can be setup to send a tagged message to syslog
    \item Rsyslog can send tags to a separate file.
    \item See lmod/contrib/tracking\_module\_usage for details
  \end{itemize}
\end{frame}

%page 44
\begin{frame}[fragile]
    \frametitle{Usage counts}
  {\small
    \begin{semiverbatim}
\$ analyzeLmodDB --sqlPattern '%fftw%' --start '2015-01-01' --end '2015-02-01'  counts

  Module path                                Distinct Users
  -----------                                --------------  
  /apps/intel13/mpich\_3\_2/mfiles/fftw3/3.3.2             151
  /apps/intel13/mpich\_3\_2/mfiles/fftw2/2.1.5              62
  /apps/intel13/impi\_4\_1/mfiles/fftw3/3.3.2               45
  /apps/intel13/impi\_4\_1/mfiles/fftw2/2.1.5               19

    \end{semiverbatim}
}
\end{frame}

%page 45
\begin{frame}[fragile]
    \frametitle{Distinct Users}
  {\small
    \begin{semiverbatim}
\$ ./analyzeLmodDB --sqlPattern '%/settarg/%' usernames

Module path                   User Name
-----------                   ---------
/apps/mfiles/settarg/5.8      user1
/apps/mfiles/settarg/5.8      user2
/apps/mfiles/settarg/5.8      user3
/apps/mfiles/settarg/5.8.1    mclay
/apps/mfiles/settarg/5.9.1    user5
    \end{semiverbatim}
}
\end{frame}

%page 46
\begin{frame}{Why does Lmod work at all?}
  \begin{itemize}
    \item The Environment is inherited from the parent process
    \item Changes in the child's environment DOES NOT affect the parent's
    \item So how could Lmod work at all? 
  \end{itemize}
\end{frame}

%page 47
\begin{frame}{The trick is}
  \begin{itemize}
    \item The \texttt{lmod} program generates text.
    \item The module command eval's that text.
    \item \texttt{module() \{ eval \$( \$LMOD\_CMD bash "\$@"; )\}}
  \end{itemize}
\end{frame}

%page 48
\begin{frame}{Why is this important?}
  \begin{itemize}
    \item It's a useful trick to know
    \item Debugging Modulefiles:
    \item \texttt{\$LMOD\_CMD bash load} \emph{module} \texttt{2$>$
        /dev/null $>$ stdout.txt}
  \end{itemize}
\end{frame}

%page 49
\begin{frame}{Debugging Lmod}
  \begin{itemize}
    \item \texttt{module --config} : reports Lmod configuration
    \item \texttt{module -D load foo $>$ load.log}
  \end{itemize}
\end{frame}

%page 50
\begin{frame}{Tracing Lmod}
  \begin{itemize}
    \item A new feature of Lmod 7.4.4+
    \item module -T ...
    \item export LMOD\_TRACING=yes
    \item Can trace loads and how restores work.
  \end{itemize}
\end{frame}

%page 51
\begin{frame}[fragile]
    \frametitle{Tracing Example}
  {\small
    \begin{semiverbatim}
running: module --initial\_load restore
  Using collection:      /home/user/.lmod.d/default
  Setting MODULEPATH to: /apps/mfiles/Darwin:/apps/mfiles/Core
  Loading: unix (fn: /apps/mfiles/Core/unix/unix.lua)
  Loading: gcc (fn: /apps/mfiles/Darwin/gcc/5.2.lua)
  Loading: StdEnv (fn: /apps/mfiles/Core/StdEnv.lua)
    \end{semiverbatim}
}
\end{frame}

%page 52
\begin{frame}{Why tracing is important}
  \begin{itemize}
    \item A single group with similar module commands in \textasciitilde/.bashrc
    \item Yet had different modules loaded. Why?
    \item Answer: One user had a default collection which changed MODULEPATH.
  \end{itemize}
\end{frame}

%page 53
\begin{frame}{How to trace Lmod startup behavior}
  \begin{itemize}
    \item How do you trace module commands in /etc/profile.d/*.sh?
    \item Could modify /etc/profile.d/z01\_StdEnv.sh to turn on
      LMOD\_TRACING for a particular user. UGH!
    \item Install SHELL STARTUP DEBUG package instead
  \end{itemize}
\end{frame}

%page 54
\begin{frame}{SHELL STARTUP DEBUG}
  \begin{itemize}
    \item Install from shellstartupdebug.sf.net
    \item Tracks startup behavior of /etc/profile.d/*.sh
    \item Allows for the setting of env. vars before
      /etc/profile.d/*.sh is sourced.
    \item Requires modifying /etc/bashrc ...
  \end{itemize}
\end{frame}

%page 55
\begin{frame}[fragile]
    \frametitle{SHELL STARTUP DEBUG (II)}
  \begin{itemize}
    \item export SHELL\_STARTUP\_DEBUG=1 (in \textasciitilde/.init.sh)
    \item Example output:
  {\small
    \begin{semiverbatim}
    /etc/profile\{
      /etc/profile.d/z00\_lmod.sh\{
      \} Time = 0.0770
      /etc/profile.d/z01\_StdEnv.sh\{
      \} Time = 0.2067
      /etc/bash.bashrc\{
      \} Time = 0.2338
    \} Time = 0.3156
    \end{semiverbatim}
}
  \end{itemize}
\end{frame}

%page 56
\begin{frame}{Examples}
  \begin{itemize}
    \item Root \textbf{must} login when Lustre is not working
    \item No Lustre directory can be Root's PATH
    \item Find it with SHELL STARTUP DEBUG
    \item Tracking MPI shell startup behavior with SHELL STARTUP DEBUG
  \end{itemize}
\end{frame}

%page 57
\begin{frame}{SHELL STARTUP DEBUG and LMOD\_TRACING}
  \begin{itemize}
    \item Put export LMOD\_TRACING=yes in \textasciitilde/.init.sh
    \item Track Lmod startup issues
  \end{itemize}
\end{frame}


%page 58
\begin{frame}[fragile]
    \frametitle{Tracing Example}
  {\small
    \begin{semiverbatim}
running: module --initial\_load restore
  Using collection:      /home/user/.lmod.d/default
  Setting MODULEPATH to: /apps/mfiles/Darwin:/apps/mfiles/Core
  Loading: unix (fn: /apps/mfiles/Core/unix/unix.lua)
  Loading: gcc (fn: /apps/mfiles/Darwin/gcc/5.2.lua)
  Loading: StdEnv (fn: /apps/mfiles/Core/StdEnv.lua)
    \end{semiverbatim}
}
\end{frame}

%page 59
\begin{frame}{Advanced Topics}
  \begin{itemize}
    \item ml 
    \item Spider Cache
    \item Semantic Versioning
    \item Modulefiles can have properties
    \item Modify functions for load() and prereq()
    \item New functions pushenv()
    \item SitePackage.lua and hooks
  \end{itemize}
\end{frame}

%page 60
\begin{frame}{ml}
  \begin{itemize}
    \item I have trouble typing ``module''
    \item Needed a shortcut
    \item ml was born
    \item ml $\rightarrow$ module list
    \item ml \emph{name} $\rightarrow$ module load \emph{name}
    \item ml -a b -c d $\rightarrow$ module unload a c; module load b d; 
    \item ml av $\rightarrow$ module avail
    \item ml spider $\rightarrow$ module spider
    \item ml load spider $\rightarrow$ module load spider
  \end{itemize}
\end{frame}

%page 61 
\begin{frame}{Spider Cache}
  \begin{itemize}
    \item Lmod has to know what modulefiles are in MODULEPATH
    \item It walks the directories in MODULEPATH every time!
    \item Or you can have system spider cache to speed things up
    \item See
      https://lmod.readthedocs.io/en/latest/130\_spider\_cache.html for
      details.
  \end{itemize}
\end{frame}

%page 62
\begin{frame}{Spider Cache (II)}
  \begin{itemize}
    \item It is much much faster to read a single big file than
      walking directories
    \item The spider cache \textbf{must} be up-to-date!
    \item Lmod trusts the cache if it exists!
    \item module --ignore\_cache avail to override!
  \end{itemize}
\end{frame}

%page 63
\begin{frame}{Spider Cache (III)}
  \begin{itemize}
    \item Normally only module avail and spider use the cache
    \item module -t avail never uses the cache
    \item module load \emph{name} does not use the cache
    \item Sites can configure Lmod to use the cache for loads
    \item Sites must use extreme care to keep the cache up-to-date!
  \end{itemize}
\end{frame}

%page 64
\begin{frame}{Spider Cache Advantages}
  \begin{itemize}
    \item The spider cache speeds up avail and spider greatly.
    \item All system modulefiles have been read, properties determined.
    \item Lua is quite fast at reading and interpreting a single file.
    \item This is preferable to walking the directory tree and reading
      every module.
    \item Properties are the reason why the cache was created.
    \item Properties must be read the modulefile or from the cache.
  \end{itemize}
\end{frame}

%page 65
\begin{frame}{Spider Cache Disadvantages}
  \begin{itemize}
    \item There is only one: Keeping it up-to-date.
    \item If Lmod sees a valid cache file it assumes it is correct.
    \item Otherwise what's the point.
    \item Currently loads bypass cache but avail and spider depend on it.
    \item Personal modules are not effected by system cache foo.
  \end{itemize}
\end{frame}

%page 66
\begin{frame}{Module version sorting}
  \begin{itemize}
    \item Old way lexicographically sort: 5.6 is newer that 5.10
    \item New way 5.10 is newer than 5.6
    \item Old to new: 2.4dev1, 2.4a1, 2.4rc2, 2.4, 2.4-1, 2.4.1
    \item Same: 2.4-1, 2.4p1, 2.4-p1
    \item Old to new: 3.2-shared, 3.2
  \end{itemize}
\end{frame}

%page 67
\begin{frame}{Load and prereq modify functions}
  \begin{itemize}
    \item \texttt{load(atleast("A","1.4"))}
    \item \texttt{load(between("B","1.4","2.7"))}
    \item \texttt{load(latest("C"))}
    \item \texttt{prereq(atleast("D","1.4"))}
    \item \texttt{prereq(between("E","1.4","2.7"))}
    \item \texttt{prereq(latest("F"))}
    \item Typically marked default or the latest.
  \end{itemize}
\end{frame}

%page 68
\begin{frame}[fragile]
    \frametitle{Module Properties}
  \begin{itemize}
    \item TACC is deploying Stampede with MIC accelerators
    \item Some modules will be ``MIC'' aware: mkl, fftw3, phdf5, ...
    \item Lmod will decorate these modules:
  {\tiny
    \begin{semiverbatim}
  1) unix/unix     3) ddt/ddt       5) mpich2/1.5    7) phdf5/1.8.9 {\color{blue}(m)}
  2) intel/13.0    4) mkl/mkl {\color{red}(*)}   6) petsc/3.2     8) StdEnv

  Where:
   {\color{blue}(m)}:  module is build natively for MIC
   {\color{red}(*)}:  module is build natively for MIC and offload to the MIC

   ------
   add_property("arch","mic")              -- \textgreater phdf5
   add_property("arch","mic:offload")      -- \textgreater mkl
    \end{semiverbatim}
}
  \item What properties would you like to support?
  \end{itemize}
\end{frame}


%page 69
\begin{frame}[fragile]
    \frametitle{Module Properties (II): Sticky}
  \begin{itemize}
    \item A module can be sticky.
    \item It requires ``\texttt{--force}'' to unload or purge.
    {\small
\begin{semiverbatim}
    add\_property("lmod","sticky")
\end{semiverbatim}
}
  \end{itemize}
\end{frame}

%page 70
\begin{frame}[fragile]
    \frametitle{pushenv()}
  \begin{itemize}
    \item Suppose you'd like to set \texttt{CC} in the environment.
    \item \texttt{setenv()} won't work.
    \item \texttt{pushenv()} will!
    {\small
\begin{semiverbatim}
                         #    setenv()     pushenv()
  $ module load   gcc;   # -\textgreater CC=gcc       CC=gcc
  $ module load   mpich; # -\textgreater CC=mpicc     CC=mpicc
  $ module unload mpich; # -\textgreater CC is unset  CC=gcc
  $ module unload gcc;   # -\textgreater CC is unset  CC is unset
\end{semiverbatim}
}
    \item \texttt{pushenv} keeps a private env vars: \texttt{\_\_LMOD\_STACK\_}\emph{NAME}
  \end{itemize}
\end{frame}

%page 71
\begin{frame}[fragile]
    \frametitle{always\_load()}
  \begin{itemize}
    \item The load() function is reversed on unload.
    \item The always\_load() function is a no-op on unload.
    \item Still you may want:
    {\small
\begin{semiverbatim}
  if (not isloaded("FOO")) then always\_load("FOO") end
\end{semiverbatim}
}
    \item Which is only slightly better than:
    {\small
\begin{verbatim}
  if (not isloaded("FOO") and mode()== "load") then
     load("FOO")
  end
\end{verbatim}
}
  \end{itemize}
\end{frame}

%page 72
\begin{frame}{SitePackage.lua}
  \begin{itemize}
    \item Functions for all modulefiles that a site needs
    \item Support for hooks like load\_hook()
    \item See contrib/SitePackage for examples
  \end{itemize}
\end{frame}

%page 73
\begin{frame}{Lmod can be configured for a site's need}
  \begin{itemize}
    \item Configuration allows for many things: redirected output,
      Site Name, System name, ...
    \item Over 40 different configuration variables.
    \item Support for shared home directories
    \item 14 different hooks to change the behavior (SitePackage.lua)
    \item Multiple language support: en, es, de, fr, zh, site
  \end{itemize}
\end{frame}

%page 74
\begin{frame}{Lessons Learned}
  \begin{itemize}
    \item A very small percentage of Lmod users join the mailing list.
    \item I'm able to support many requests for features but not all.
    \item Hardening:
      \begin{itemize}
        \item sandbox for evaluating modulefiles
        \item Checking argument types
        \item checkModuleSyntax script
      \end{itemize}
    \item Users sometimes ask the right question.
    \item Lmod is much better product because of its wider use.
  \end{itemize}
\end{frame}

%page 75
\begin{frame}{Lessons Learned (II)}
  \begin{itemize}
    \item It is hard work to develop a communities trust.
    \item Some site will be managed in a way that I have never dreamed of.
    \item Deprecating features is difficult with Lmod.
    \item Users would get deprecated feature reports that sys-admins
      must fix.
  \end{itemize}
\end{frame}


%page 76
\begin{frame}{Regression Testing of Lmod}
  \begin{itemize}
    \item A suite of 100+ tests each with many steps.
    \item No release without passing all those tests.
    \item These tests make Lmod re-factoring much easier.
    \item The github repo is generally safe.
  \end{itemize}
\end{frame}

%page 77
\begin{frame}{Remote Debugging}
  \begin{itemize}
    \item No software over ten lines is bug free.
    \item Lmod is no exception.
    \item Bug reports are as easy as:
      \begin{itemize}
        \item \texttt{module --config  2\textgreater  config.log}
        \item \texttt{module -D avail  2\textgreater  avail.log}
      \end{itemize}
    \item Sometimes I'll create test versions with more debugging for
      you to test.
  \end{itemize}
\end{frame}

%page 78
\begin{frame}{Recommendations for Good Site Management}
  \begin{itemize}
    \item A module purge should not break things.
    \item Set \texttt{LMOD\_SYSTEM\_DEFAULT\_MODULES} so that module
      restore and reset works.
    \item Encourage the use of "module save" so that users can have a
      default set of modules.
    \item Use the family directive to protect users.
    \item Consider the use of "checkModuleSyntax" before installing new modules.
  \end{itemize}
\end{frame}

%page 79
\begin{frame}{Conclusions: Lmod}
  \begin{itemize}
    \item Latest version: https://github.com:TACC/Lmod.git
    \item Stable version: http://lmod.sf.net
    \item Documentation:  http://lmod.readthedocs.org
    \item Shell Startup Debug: http://shellstartupdebug.sf.net
    \item Mailing List:   lmod-users@lists.sourceforge.net.
    \item Join here: https://lists.sourceforge.net/lists/listinfo/lmod-users
  \end{itemize}
\end{frame}


%
\begin{frame}{License}
\footnotesize
\textcopyright The University of Texas at Austin, \the\year\\
\vspace{0.5 cm}
This work is licensed under the Creative Commons Attribution Non-Commercial 3.0 Unported License. To view a copy of this license, 
visit http://creativecommons.org/licenses/by-nc/3.0/\\
\vspace{0.5 cm}
When attributing this work, please use the following text: 
"Designing and Administering Large-scale Systems", Texas Advanced Computing Center, \the\year. Available under a Creative Commons Attribution Non-Commercial 3.0 Unported License.
\newline
\newline
\includegraphics[width=1cm]{./figures/cc.png}
\end{frame}

\end{document}
